\chapter{IMPLEMENTATION}
\label{ch:implementation}


\hs This chapter presents the software implementation of the \textbf{kfc\_procedure} Python package. The package contains two connected machine learning subsystems: the \KFCProcedure{} clusterwise learning framework and the \GradientCOBRA{} prediction-space aggregation framework. These subsystems are implemented within a single Python package so that clusterwise local learning and kernel-based aggregation can be used independently or combined within one workflow.


Building on the methodological design introduced in \chapref{ch:methodology}, this chapter focuses on how the proposed framework is realized in code. The discussion covers repository organization, software design principles, interface design, factory-based component registration, clustering and divergence implementations, COBRA aggregation components, orchestration logic, utility functions, testing, documentation, and packaging infrastructure.

The implementation follows the structure of a research-oriented Python library rather than a large-scale application architecture. It uses a \textbf{src/} package layout, scikit-learn-style estimator interfaces, modular factory registries, and component-level tests for the COBRA core subsystem. Therefore, this chapter emphasizes the traceability between the repository structure and the implemented machine learning workflows.

% ============================
\section{Source Code Architecture}
\label{sec:repository_architecture}

\hs The repository is organized as a Python package using a \textbf{src/} layout. This structure separates the importable package implementation from experiment notebooks, tests, package configuration files, documentation assets, and automation workflows. Figure~\ref{fig:kfcprocedure_repository_structure} shows the overall repository structure.

\subsection{Top-Level Repository Structure}

\begin{figure}[H]
      \centering
      \begin{minipage}{0.85\textwidth}
            \begin{Verbatim}[fontsize=\footnotesize, baselinestretch=0.9]
                  kfc-procedure/
                  ├── .github/
                  │   └── workflows/
                  │       └── ci-cd.yml
                  ├── notebooks/
                  │   ├── kfc.ipynb
                  │   └── experiments/
                  │       ├── classification/
                  │       └── regression/
                  ├── src/
                  │   └── kfc_procedure/
                  ├── tests/
                  │   └── cobra/
                  ├── pyproject.toml
                  ├── requirements.txt
                  ├── README.md
                  ├── LICENSE
                  └── .pylintrc
            \end{Verbatim}
      \end{minipage}
      \caption{Overall repository structure of the KFCProcedure project}
      \label{fig:kfcprocedure_repository_structure}
\end{figure}

\hs As shown in Figure~\ref{fig:kfcprocedure_repository_structure}, the repository separates source code, notebooks, tests, packaging configuration, and automation files. The package configuration defines the package metadata, Python version requirement, dependencies, optional extras, build configuration, and pytest settings. The distribution package is named \textbf{kfc-procedure}, while the importable Python package is named \textbf{kfc\_procedure}. The project requires Python~3.11 or newer.

\subsection{Main Package Structure}

\hs The implementation package is organized into a high-level orchestration module, a KFC core subsystem, a COBRA aggregation subsystem, and utility components. Figure~\ref{fig:kfcprocedure_package_structure} shows the main internal package structure.

\begin{figure}[H]
      \centering
      \begin{minipage}{0.85\textwidth}
            \begin{Verbatim}[fontsize=\footnotesize, baselinestretch=0.9]
                  kfc_procedure/
                  ├── kfc.py
                  ├── core/
                  │   ├── clustering/
                  │   ├── combiner/
                  │   ├── ml/
                  │   ├── steps/
                  │   └── factory.py
                  ├── cobra/
                  │   ├── gradientcobra.py
                  │   ├── mixcobra.py
                  │   ├── combined_classifier.py
                  │   ├── core/
                  │   └── utils/
                  └── utils/
                  ├── logger.py
                  └── resolve.py
            \end{Verbatim}
      \end{minipage}
      \caption{Main package structure of the KFCProcedure framework}
      \label{fig:kfcprocedure_package_structure}
\end{figure}

As shown in Figure~\ref{fig:kfcprocedure_package_structure}, the framework is separated into four main implementation layers. Table~\ref{tab:implementation_layers} summarizes the role of each layer.

\begin{table}[H]
      \centering
      \caption{Main Implementation Layers of the Framework}
      \label{tab:implementation_layers}
      \renewcommand{\arraystretch}{1.2}
      \begin{tabularx}{\textwidth}{
                  >{\raggedright\arraybackslash}p{3.5cm}
                  >{\raggedright\arraybackslash}X
            }
            \toprule
            \textbf{Layer}          & \textbf{Description}                                                                                                                                \\
            \midrule
            KFC core layer          &
            Provides the main clusterwise learning components, including Bregman clustering, local model wrappers, K-step/F-step/C-step workflow components, and KFC combiner strategies. \\

            COBRA aggregation layer &
            Provides COBRA-based aggregation models, including \textbf{GradientCOBRA}, \textbf{MixCOBRARegressor}, \textbf{CombinedClassifier}, and reusable aggregation components.      \\

            Orchestration layer     &
            Exposes the high-level classes \textbf{KFCProcedure}, \textbf{KFCRegressor}, and \textbf{KFCClassifier}, which coordinate training and prediction workflows.                  \\

            Utility layer           &
            Provides supporting functions for logging, component resolution, data splitting, normalization constants, estimator management, and distance-related helper operations.       \\
            \bottomrule
      \end{tabularx}


\end{table}

\subsection{Clusterwise Learning Subsystem}

\hs The clusterwise learning subsystem forms the KFC core of the framework. It is responsible for constructing divergence-specific clusters, training local predictive models inside those clusters, and combining the resulting prediction outputs. Table~\ref{tab:kfc_core_components} summarizes the main components of this subsystem.

\begin{table}[H]
      \centering
      \caption{Main Components of the Clusterwise Learning Subsystem}
      \label{tab:kfc_core_components}
      \renewcommand{\arraystretch}{1.2}
      \begin{tabularx}{\textwidth}{
                  >{\raggedright\arraybackslash}p{3.5cm}
                  >{\raggedright\arraybackslash}X
            }
            \toprule
            \textbf{Component} & \textbf{Role}                                                                    \\
            \midrule
            Clustering         &
            Implements divergence-based clustering using Bregman K-Means and registered divergence functions.     \\

            Local learning     &
            Provides local predictive models and adapters for training cluster-specific models during the F-step. \\

            Workflow step      &
            Defines the K-step, F-step, and C-step execution logic used in the clusterwise learning pipeline.     \\

            Combiner           &
            Implements final aggregation strategies for combining divergence-level prediction outputs.            \\

            Factory            &
            Resolves divergences, local models, and combiners through registered component names.                 \\
            \bottomrule
      \end{tabularx}

\end{table}

The clusterwise learning subsystem therefore follows the structure of the proposed KFCProcedure. The K-step constructs divergence-specific cluster assignments, the F-step trains local models within the clusters, and the C-step combines the divergence-level prediction matrix into a final prediction.

\subsection{COBRA Aggregation Subsystem}

\hs The COBRA aggregation subsystem provides models and reusable components for prediction-space aggregation. It supports regression aggregation through \textbf{GradientCOBRA} and \textbf{MixCOBRARegressor}, classification aggregation through \textbf{CombinedClassifier}. Table~\ref{tab:cobra_subsystem_components} summarizes the main components of the COBRA subsystem.

\begin{table}[H]
      \centering
      \caption{Main Components of the COBRA Aggregation Subsystem}
      \label{tab:cobra_subsystem_components}
      \renewcommand{\arraystretch}{1.2}
      \begin{tabularx}{\textwidth}{
                  >{\raggedright\arraybackslash}p{3.8cm}
                  >{\raggedright\arraybackslash}X
            }
            \toprule
            \textbf{Component}        & \textbf{Role}                                                                                                                  \\
            \midrule
            GradientCOBRA             &
            Provides regression aggregation based on prediction-space distances and kernel-weighted calibration targets.                                               \\

            MixCOBRARegressor         &
            Combines input-space and prediction-space distance information for regression aggregation.                                                                 \\

            CombinedClassifier        &
            Provides classification aggregation by combining the predictions of multiple classification models.                                                        \\

            Reusable COBRA components &
            Provide shared infrastructure for distances, kernels, losses, optimizers, splitters, normalizers, adapters, cross-validation, estimators, and aggregators. \\
            \bottomrule
      \end{tabularx}

\end{table}

The main COBRA classes reuse the same component infrastructure while implementing different aggregation workflows. This design avoids duplicating common logic for distance computation, kernel weighting, optimization, and aggregation.

\subsection{Architectural Summary}

\hs The codebase implements a two-subsystem architecture. The KFC subsystem builds divergence-specific clusterwise predictors, while the COBRA subsystem aggregates predictions using distance-based and kernel-based mechanisms. The two subsystems interact through the KFC C-step, where COBRA-based combiners can be selected as aggregation strategies. This structure supports modular experimentation because clustering, local learning, prediction aggregation, optimization, and utility operations are separated into specialized components.


\section{Software Design Principles}
\label{sec:software_design_principles}

\hs The \textbf{\textit{kfc\_procedure}} framework is designed according to a set of software engineering principles that support modularity, extensibility, maintainability, and interoperability. Since the framework combines clusterwise supervised learning with COBRA-based prediction-space aggregation, the implementation requires clear structural boundaries between orchestration logic, clustering procedures, local learning models, aggregation mechanisms, optimization components, and supporting utilities.

This section presents the main design principles applied throughout the implementation. Each principle is described in relation to the repository structure and illustrated using architectural diagrams to show how the design is reflected in the actual framework.

% ============================================================
\subsection{Layered and Modular Architecture}

\hs The framework adopts a layered and modular architecture in which high-level orchestration logic is separated from lower-level algorithmic components. This organization allows the main workflow coordination to remain independent from the implementation details of clustering, local model training, kernel computation, optimization, and aggregation. As a result, the system becomes easier to maintain, test, and extend.

\begin{figure}[H]
      \centering
      \begin{tikzpicture}[
                  node distance=1.4cm,
                  box/.style={
                              draw,
                              rounded corners,
                              align=center,
                              text width=4.5cm,
                              font=\small
                        },
                  arrow/.style={-Latex, thick}
            ]

            \node[box] (top) {Orchestration Layer\\ \textbf{kfc}};

            \node[box, below left=of top] (kfc) {Clusterwise Learning\\ \textbf{core/}};
            \node[box, below right=of top] (cobra) {COBRA Aggregation\\ \textbf{cobra/}};

            \node[box, below=of kfc] (kfc2) {
                  Clustering,\\ ML Models, \\and Pipeline Steps
            };

            \node[box, below=of cobra] (cobra2) {
                  Kernels,\\ Optimizers, \\Distances, Aggregators
            };

            \draw[arrow] (top) -- (kfc);
            \draw[arrow] (top) -- (cobra);

            \draw[arrow] (kfc) -- (kfc2);
            \draw[arrow] (cobra) -- (cobra2);

      \end{tikzpicture}
      \caption{Layered modular architecture of the \textbf{kfc\_procedure} framework}
      \label{fig:layered-modular-architecture}
\end{figure}

As illustrated in \figref{fig:layered-modular-architecture}, the orchestration layer implemented in \textbf{kfc.py} provides the high-level entry point of the framework. This layer coordinates the execution of the K-step, F-step, and C-step, while the lower-level packages provide the algorithmic components required by each stage. The \textbf{core/} package contains the clusterwise learning components, whereas the \textbf{cobra/} package contains the COBRA-based aggregation components.

This architectural separation reduces coupling between the main workflow and the underlying algorithms. Consequently, clustering methods, learning models, and aggregation mechanisms can be modified or extended without requiring major changes to the orchestration layer.

% ============================================================

\subsection{Interface-Driven Design}

\hs The framework applies interface-driven design by defining common contracts for major component families. Instead of requiring the pipeline to depend directly on concrete implementations, components are accessed through shared interfaces or scikit-learn-style method conventions such as \textbf{fit()}, \textbf{predict()}, and, where applicable, \textbf{predict\_proba()}. This design improves consistency across different modules and allows alternative implementations to be integrated more easily.

\begin{figure}[H]
      \centering
      \begin{tikzpicture}[
                  node distance=1.6cm,
                  box/.style={
                              draw,
                              rounded corners,
                              align=center,
                              text width=3.8cm,
                              font=\small,
                              minimum height=1.1cm
                        },
                  arrow/.style={-Latex, thick}
            ]

            \node[box] (base) {Abstract Component Interface};

            \node[box, below left=of base] (a) {Local Model\\Interface};
            \node[box, below=of base] (b) {Kernel / Distance\\Interface};
            \node[box, below right=of base] (c) {Optimizer / Combiner\\Interface};

            \draw[arrow] (base) -- (a);
            \draw[arrow] (base) -- (b);
            \draw[arrow] (base) -- (c);

      \end{tikzpicture}
      \caption{Interface-driven design across reusable framework components}
      \label{fig:interface-driven-design}
\end{figure}

As shown in \figref{fig:interface-driven-design}, the framework uses abstract component families to organize local models, kernels, distances, optimizers, and combiners. For example, local learning models are required to expose common training and prediction behavior, while COBRA components such as kernels, distances, losses, and optimizers follow their own reusable interfaces.

This design allows the framework to support multiple interchangeable implementations while preserving a stable pipeline structure. Therefore, new estimators, kernels, optimizers, or aggregation methods can be introduced without changing the high-level training and prediction workflow.

% ============================================================

\subsection{Composition-Based Design}

\hs The implementation follows a composition-based design, where complex workflows are constructed by combining independent components rather than relying on deep inheritance hierarchies. In the KFC pipeline, the main procedure is composed of separate K-step, F-step, and C-step objects, each of which encapsulates a specific part of the learning process.

\begin{figure}[H]
      \centering
      \begin{tikzpicture}[
                  node distance=1.4cm,
                  box/.style={
                              draw,
                              rounded corners,
                              align=center,
                              text width=4.2cm,
                              font=\small,
                              minimum height=1.1cm
                        },
                  arrow/.style={-Latex, thick},
                  dashedarrow/.style={-Latex, thick, dashed}
            ]

            \node[box] (kfc) {KFCProcedure};

            \node[box, below left=of kfc] (kstep) {KStep\\Divergence-Based Clustering};
            \node[box, below=of kfc] (fstep) {FStep\\Local Models per Cluster};
            \node[box, below right=of kfc] (cstep) {CStep\\Combiner / Aggregation};

            \node[box, below=1.4cm of kstep] (div) {Bregman Divergences\\Multiple Metrics};
            \node[box, below=1.4cm of fstep] (est) {Local Estimators\\Cluster-Level Learners};
            \node[box, below=1.4cm of cstep] (comb) {Aggregation Strategy\\Voting / Weighted / COBRA};

            \node[box, left=1.5cm of kfc] (sk) {scikit-learn-style\\Estimator API};

            \draw[arrow] (kfc) -- (kstep);
            \draw[arrow] (kfc) -- (fstep);
            \draw[arrow] (kfc) -- (cstep);

            \draw[arrow] (kstep) -- (div);
            \draw[arrow] (fstep) -- (est);
            \draw[arrow] (cstep) -- (comb);

            \draw[dashedarrow] (kfc) -- (sk);

      \end{tikzpicture}
      \caption{Composition-based organization of the KFC pipeline}
      \label{fig:composition-based-design}
\end{figure}

As illustrated in \figref{fig:composition-based-design}, \textbf{KFCProcedure} does not directly implement all clustering, learning, and aggregation logic inside one monolithic class. Instead, it coordinates independent pipeline components. The K-step handles divergence-based clustering, the F-step trains local models within discovered clusters, and the C-step combines the divergence-level prediction outputs.

This design improves maintainability because each pipeline stage can be developed and tested independently. It also improves extensibility because new clustering divergences, local estimators, or aggregation strategies can be composed into the existing workflow with limited changes to the overall architecture.

% ============================================================

\subsection{Strategy Design Pattern}

\hs The framework applies the strategy design pattern by implementing alternative algorithms as interchangeable strategies. This principle is especially important in the aggregation stage, where different prediction combination methods can be selected depending on the task, model configuration, or experimental objective.

\begin{figure}[H]
      \centering
      \begin{tikzpicture}[
                  node distance=1.4cm,
                  box/.style={
                              draw,
                              rounded corners,
                              align=center,
                              text width=3.8cm,
                              font=\small,
                              minimum height=1.1cm
                        },
                  arrow/.style={-Latex, thick}
            ]

            \node[box] (s) {Aggregation Strategy};

            \node[box, below left=of s] (m) {Weighted Mean};
            \node[box, below=of s] (v) {Weighted Vote};
            \node[box, below right=of s] (g) {GradientCOBRA};

            \draw[arrow] (s) -- (m);
            \draw[arrow] (s) -- (v);
            \draw[arrow] (s) -- (g);

      \end{tikzpicture}
      \caption{Strategy pattern for interchangeable aggregation methods}
      \label{fig:strategy-design-pattern}
\end{figure}

As shown in \figref{fig:strategy-design-pattern}, the aggregation stage can use different strategies, including rule-based aggregation, weighted voting, weighted averaging, stacking-based methods, or COBRA-based aggregation. These strategies share the same role within the pipeline but differ in their internal prediction-combination mechanisms.

This design enables runtime selection of aggregation methods without modifying the core KFC workflow. It also supports fair experimentation, since different aggregation strategies can be compared under the same pipeline structure.

% ============================================================

\subsection{Factory-Based Construction}

\hs The framework adopts factory-based construction to support dynamic component creation from configuration parameters. Instead of directly instantiating specific classes inside the pipeline, the system delegates component creation to factory and registry mechanisms. This approach separates object construction from workflow execution.

\begin{figure}[H]
      \centering
      \begin{tikzpicture}[
                  node distance=1.4cm,
                  box/.style={
                              draw,
                              rounded corners,
                              align=center,
                              text width=4.0cm,
                              font=\small,
                              minimum height=1.1cm
                        },
                  arrow/.style={-Latex, thick}
            ]

            \node[box] (cfg) {Configuration Input};

            \node[box, below=of cfg] (f) {Factory and Registry Layer\\ \textbf{factory.py}};

            \node[box, below left=of f] (c1) {Clusterer / Divergence};
            \node[box, below=of f] (c2) {Estimator / Combiner};
            \node[box, below right=of f] (c3) {Kernel / Optimizer};

            \draw[arrow] (cfg) -- (f);
            \draw[arrow] (f) -- (c1);
            \draw[arrow] (f) -- (c2);
            \draw[arrow] (f) -- (c3);

      \end{tikzpicture}
      \caption{Factory-based construction of configurable framework components}
      \label{fig:factory-based-construction}
\end{figure}

As illustrated in \figref{fig:factory-based-construction}, the factory layer receives configuration inputs and resolves them into concrete framework components. For example, a user may select a divergence, estimator, kernel, optimizer, or combiner by name, while the factory layer is responsible for instantiating the corresponding implementation.

This mechanism improves reproducibility and extensibility. Experimental configurations can be changed without altering the pipeline source code, and new algorithms can be registered in the factory system while preserving the existing interface used by the orchestration layer.

% ==============================

\subsection{Separation of Concerns}

\hs The implementation follows the principle of separation of concerns by assigning each module a specific and clearly defined responsibility. This principle prevents unrelated functionalities from being mixed within the same component and ensures that clustering, local learning, kernel computation, optimization, aggregation, and support operations remain structurally independent.

\begin{figure}[H]
      \centering
      \begin{tikzpicture}[
                  box/.style={
                              draw,
                              rounded corners,
                              align=center,
                              text width=3.4cm,
                              font=\small,
                              minimum height=1.4cm
                        },
                  center/.style={
                              draw,
                              rounded corners,
                              align=center,
                              text width=4.5cm,
                              fill=gray!10,
                              font=\small,
                              minimum height=1.5cm
                        },
                  arrow/.style={-Latex, thick}
            ]

            \node[center] (sys) {KFCProcedure Framework};

            \node[box, above=1.5cm of sys] (a) {Clustering Module};
            \node[box, left=1.8cm of sys] (b) {Local Learning Module};
            \node[box, right=1.8cm of sys] (c) {Kernel Computation Module};
            \node[box, below left=1.4cm of sys] (d) {Optimization Module};
            \node[box, below right=1.4cm of sys] (e) {Utility Layer};

            \draw[arrow] (sys) -- (a);
            \draw[arrow] (sys) -- (b);
            \draw[arrow] (sys) -- (c);
            \draw[arrow] (sys) -- (d);
            \draw[arrow] (sys) -- (e);

      \end{tikzpicture}
      \caption{Separation of concerns in the KFCProcedure framework}
      \label{fig:separation_of_concerns}
\end{figure}

As shown in Figure~\ref{fig:separation_of_concerns}, the framework is organized around specialized modules with distinct responsibilities. The clustering module is responsible for divergence-based clustering and cluster assignment. The local learning module provides the predictive models used within the F-step. The kernel computation module supports kernel-based similarity transformations for COBRA-style aggregation, while the optimization module provides procedures for tuning aggregation-related parameters. The utility layer supplies reusable support functions such as logging, component resolution, data splitting, normalization, and other auxiliary operations.

This separation allows each module to evolve independently while maintaining a coherent overall system structure. For example, a new divergence function can be introduced into the clustering module without requiring changes to the optimization module, provided that the expected interfaces remain consistent. Therefore, separation of concerns contributes directly to the clarity, maintainability, and scalability of the implementation.

% ==========================================

\subsection{Extensibility-Oriented Design}

\hs The framework is designed to support extensibility across the three main stages of the KFC learning pipeline: K-step, F-step, and C-step. Each stage exposes stable component boundaries, allowing new divergences, local estimators, and aggregation strategies to be integrated without modifying the overall workflow.

Extensibility is achieved through a combination of abstract interfaces, factory registries, and configuration-driven component selection. This design enables the framework to support both existing implementations and future algorithmic extensions.

\begin{figure}[H]
      \centering

      \resizebox{0.98\linewidth}{!}{%
            \begin{tikzpicture}[
                        box/.style={
                                    draw,
                                    rounded corners,
                                    align=center,
                                    text width=4.3cm,
                                    minimum height=1.3cm,
                                    font=\small
                              },
                        mainbox/.style={box, fill=gray!10},
                        newbox/.style={box, fill=gray!5},
                        arrow/.style={-Latex, thick}
                  ]

                  \node[mainbox] (system) {KFCProcedure System\\
                        \footnotesize K-Step $\rightarrow$ F-Step $\rightarrow$ C-Step};

                  \node[box, below=1.2cm of system] (interface)
                  {Abstract Interfaces\\
                        \footnotesize Divergence / Estimator / Combiner};

                  \node[box, below=1.2cm of interface] (factory)
                  {Factory \& Registry Layer\\
                        \footnotesize Configuration-Driven Component Selection};

                  \node[box, below left=1.5cm and 1.2cm of factory] (existing)
                  {Existing Components\\
                        \footnotesize
                        K-Step: registered divergences\\
                        F-Step: local estimators\\
                        C-Step: registered combiners};

                  \node[newbox, below right=1.5cm and 1.2cm of factory] (new)
                  {Pluggable Extensions\\
                        \footnotesize
                        new divergences for K-Step\\
                        new estimators for F-Step\\
                        new combiners for C-Step};

                  \draw[arrow] (system) -- (interface);
                  \draw[arrow] (interface) -- (factory);
                  \draw[arrow] (factory) -- (existing);
                  \draw[arrow] (factory) -- (new);

            \end{tikzpicture}%
      }

      \caption{Extensibility of the KFCProcedure Framework}
      \label{fig:extensibility-oriented-design}
\end{figure}

As illustrated in \figref{fig:extensibility-oriented-design}, extensibility is supported at multiple levels of the pipeline. The K-step can be extended by introducing additional divergence measures for clustering. The F-step can be extended by registering new supervised learning models as local estimators. The C-step can be extended by adding new aggregation strategies, including rule-based, stacking-based, or COBRA-based combiners.

This design is particularly useful for research-oriented machine learning frameworks, where experimental flexibility is required. New methods can be added and evaluated within the same architectural structure, which improves comparability, reproducibility, and long-term maintainability.

% =============================================

\subsection{Summary}

\hs Overall, the \textbf{\textit{kfc\_procedure}} framework is implemented using a modular, layered, interface-driven, and extensibility-oriented software architecture. The combination of composition, strategy-based selection, factory-based construction, and separation of concerns ensures that the framework remains flexible and maintainable while supporting multiple machine learning workflows.

These design principles are important because the framework integrates two related but distinct subsystems: KFC clusterwise learning and COBRA prediction-space aggregation. By maintaining clear architectural boundaries between these subsystems, the implementation supports independent development, systematic experimentation, and future extension of the proposed framework.

% ============================

\section{Core Interface Layer}
\label{sec:core_interface_layer}

\hs The framework does not rely on a single universal interface layer. Instead, it defines several interface families for different parts of the system. The KFC subsystem provides abstractions for divergences, local models, and combiners, while the COBRA subsystem provides abstractions for estimators, splitters, distances, kernels, losses, optimizers, normalizers, adapters, cross-validation strategies, and aggregators.

This design keeps the two subsystems modular while still allowing them to interact through the prediction matrix produced by the KFC F-step and passed to the C-step combiner.

\subsection{Factory Infrastructure}

\hs The framework uses registry-based factory mechanisms to create and manage configurable components. These factories allow components to be selected by string identifiers instead of requiring users to instantiate each object manually. Table~\ref{tab:factory_infrastructure} summarizes the two main factory infrastructures used in the framework.

\begin{table}[H]
      \centering
      \caption{Factory Infrastructure in the Framework}
      \label{tab:factory_infrastructure}
      \renewcommand{\arraystretch}{1.2}
      \begin{tabularx}{\textwidth}{
                  >{\raggedright\arraybackslash}p{3.4cm}
                  >{\raggedright\arraybackslash}X
            }
            \toprule
            \textbf{Factory Scope}       & \textbf{Supported Components}                                                                                                                  \\
            \midrule
            KFC factory infrastructure   &
            Supports KFC-related components such as divergence factories, local model factories, and combiner factories.                                                                  \\

            COBRA factory infrastructure &
            Supports COBRA-related components such as estimators, splitters, distances, kernels, losses, optimizers, normalizers, adapters, cross-validation strategies, and aggregators. \\
            \bottomrule
      \end{tabularx}
\end{table}

Both factory systems support case-insensitive registration, aliases, category labels, runtime object creation, and discovery of available components. This allows users to select components through identifiers such as \textbf{euclidean}, \textbf{rbf}, \textbf{mse}, \textbf{grid}, or \textbf{weighted\_mean}. As a result, the framework remains configurable while preserving a consistent interface for component resolution.

\subsection{KFC Interface Families}

\hs The KFC subsystem relies on several core interface families to support divergence registration, local model training, and prediction aggregation. These abstractions provide a consistent way to extend the framework with new divergences, local learners, and combiners. Table~\ref{tab:kfc_interface_families} summarizes the main KFC interface families.

\begin{table}[H]
      \centering
      \caption{KFC Interface Families}
      \label{tab:kfc_interface_families}
      \renewcommand{\arraystretch}{1.2}
      \begin{tabularx}{\textwidth}{
                  >{\raggedright\arraybackslash}p{3.8cm}
                  >{\raggedright\arraybackslash}X
            }
            \toprule
            \textbf{Interface Family}    & \textbf{Role}                                                                                                                         \\
            \midrule
            Bregman divergence interface &
            Defines the structure for divergence functions used by \textbf{BregmanKMeans}. Concrete divergences implement the required mathematical functions and domain checks. \\

            Divergence factory           &
            Registers and instantiates supported divergence implementations, including \textbf{euclidean}, \textbf{gkl}, \textbf{is}, and \textbf{logistic}.                     \\

            Local model interface        &
            Defines the interface used by local predictive models in the F-step. Local models are expected to provide fitting and prediction behavior.                           \\

            Local model factory          &
            Registers local models and resolves them by name. It includes a native mean regressor and compatible scikit-learn classifiers and regressors.                        \\

            Combiner interface           &
            Defines the interface for C-step combiners used to aggregate divergence-level prediction outputs.                                                                    \\

            Combiner factory             &
            Registers regression and classification combiners, including simple aggregation rules, stacking-based combiners, and COBRA-based combiners.                          \\
            \bottomrule
      \end{tabularx}
\end{table}

The clustering workflow is handled directly by \textbf{BregmanKMeans} and \textbf{KStep}. The current implementation does not define a separate generic \textbf{BaseClusterer} abstraction.

\subsection{COBRA Interface Families}

\hs The COBRA subsystem defines a set of reusable interface families that support estimator management, data splitting, distance computation, kernel weighting, optimization, aggregation, cross-validation, and normalization. These interfaces follow a factory-based design so that different components can be selected and configured consistently. Table~\ref{tab:cobra_interface_families} summarizes the main COBRA interface families.

\begin{table}[H]
      \centering
      \caption{COBRA Interface Families}
      \label{tab:cobra_interface_families}
      \renewcommand{\arraystretch}{1.2}
      \begin{tabularx}{\textwidth}{
                  >{\raggedright\arraybackslash}p{5cm}
                  >{\raggedright\arraybackslash}X
            }
            \toprule
            \textbf{Interface Family}   & \textbf{Role}                                                                                   \\
            \midrule
            Data splitter interfaces    &
            Support train/calibration splitting strategies used during COBRA model fitting.                                               \\

            Estimator interfaces        &
            Define and resolve COBRA base learners through estimator abstractions and factory registration.                               \\

            Distance interfaces         &
            Provide pairwise distance computation for prediction-space or input-space comparison.                                         \\

            Kernel interfaces           &
            Transform computed distances into aggregation weights.                                                                        \\

            Kernel adapter interfaces   &
            Support one-parameter and two-parameter transformations of distance matrices before kernel weighting.                         \\

            Loss interfaces             &
            Define optimization objectives used when tuning aggregation parameters.                                                       \\

            Optimizer interfaces        &
            Provide grid-search and gradient-based methods for tuning parameters such as bandwidth, $\alpha$, and $\beta$.                \\

            Aggregator interfaces       &
            Combine calibration targets through weighted mean aggregation for regression or weighted vote aggregation for classification. \\

            Cross-validation interfaces &
            Provide K-fold, stratified K-fold, and time-series validation strategies for internal parameter selection.                    \\

            Normalizer interfaces       &
            Provide reusable standard and min--max normalization utilities for model inputs or prediction-space values.                   \\
            \bottomrule
      \end{tabularx}
\end{table}

\subsection{Scikit-Learn Compatibility}

\hs The high-level estimators inherit from or follow scikit-learn-style interfaces. \textbf{KFCProcedure}, \textbf{KFCRegressor}, and \textbf{KFCClassifier} expose \textbf{fit()} and \textbf{predict()} methods. The COBRA classes \textbf{GradientCOBRA} and \textbf{MixCOBRARegressor} inherit from scikit-learn estimator and regressor mixin classes, while \textbf{CombinedClassifier} follows a classifier-style interface with \textbf{fit()}, \textbf{predict()}, and \textbf{predict\_proba()}.

One implementation limitation should be noted. The full \textbf{KFCProcedure.predict\_proba()} method calls \textbf{FStep.predict\_proba()}, but the current \textbf{FStep} class does not implement this method. Therefore, full KFC probability prediction should not be described as a completed feature in the current implementation.

% ============================

\section{Clusterwise Learning Subsystem}
\label{sec:kfc_subsystem}

\hs The clusterwise learning subsystem forms the core KFC component of the library. It integrates divergence-based clustering, cluster-specific predictive modeling, and combiner-based prediction fusion to support supervised learning on heterogeneous data structures.

\subsection{Subsystem Architecture}

\hs The KFC subsystem is organized around five main architectural components: clustering, local learning, workflow steps, prediction combiners, and factory-based component resolution. These components separate the responsibilities of the clusterwise learning pipeline and make the framework easier to extend and maintain. Table~\ref{tab:kfc_subsystem_architecture} summarizes the main components of the KFC subsystem.

\begin{table}[H]
      \centering
      \caption{Main Components of the KFC Subsystem}
      \label{tab:kfc_subsystem_architecture}
      \renewcommand{\arraystretch}{1.2}
      \begin{tabularx}{\textwidth}{
                  >{\raggedright\arraybackslash}p{3.5cm}
                  >{\raggedright\arraybackslash}X
            }
            \toprule
            \textbf{Component}   & \textbf{Role in the KFC Subsystem}                                         \\
            \midrule
            Clustering           &
            Provides divergence-based clustering through Bregman K-Means and registered divergence functions. \\

            Local learning       &
            Provides local predictive models used within the clusters produced by the K-step.                 \\

            Workflow steps       &
            Defines the K-step, F-step, and C-step execution logic of the clusterwise learning pipeline.      \\

            Prediction combiners &
            Combines the divergence-level prediction outputs produced by the F-step.                          \\

            Factory resolution   &
            Resolves clustering models, local models, and combiners from registered component names.          \\
            \bottomrule
      \end{tabularx}
\end{table}

The implementation separates the clusterwise workflow into three main pipeline stages, as summarized in Table~\ref{tab:kfc_pipeline_stages}.

\begin{table}[H]
      \centering
      \caption{Main Pipeline Stages of the KFC Workflow}
      \label{tab:kfc_pipeline_stages}
      \renewcommand{\arraystretch}{1.2}
      \begin{tabularx}{\textwidth}{
                  >{\raggedright\arraybackslash}p{2.5cm}
                  >{\raggedright\arraybackslash}X
            }
            \toprule
            \textbf{Stage} & \textbf{Description}                                                                             \\
            \midrule
            K-step         &
            Fits one Bregman K-Means model for each selected divergence and produces divergence-specific cluster assignments. \\

            F-step         &
            Trains local predictive models for each divergence and observed cluster.                                          \\

            C-step         &
            Combines the divergence-level prediction matrix using the selected registered combiner.                           \\
            \bottomrule
      \end{tabularx}
\end{table}

\subsection{Clustering Module}

\hs The clustering module provides the divergence-based clustering mechanism used in the K-step of the \KFCProcedure{} framework. Its main role is to partition the input data into clusters under different divergence assumptions before local predictive models are trained within each cluster.

The central class of this module is \textbf{BregmanKMeans}. It follows a Lloyd-style clustering procedure and provides several implementation features that support repeated initialization, divergence-based sample assignment, centroid updating, and scikit-learn-style usage. Table~\ref{tab:bregman_kmeans_features} summarizes the main features of this clustering component.

\begin{table}[H]
      \centering
      \caption{Main Features of BregmanKMeans}
      \label{tab:bregman_kmeans_features}
      \renewcommand{\arraystretch}{1.2}
      \begin{tabularx}{\textwidth}{
                  >{\raggedright\arraybackslash}p{4.2cm}
                  >{\raggedright\arraybackslash}X
            }
            \toprule
            \textbf{Feature}             & \textbf{Description}                                                                                                              \\
            \midrule
            Centroid initialization      &
            Initializes centroids randomly from training samples.                                                                                                            \\

            Multiple initializations     &
            Supports repeated random initializations through \varname{n\_init}.                                                                                              \\

            Divergence-based assignment  &
            Assigns samples to clusters using the selected Bregman divergence.                                                                                               \\

            Centroid update              &
            Updates cluster centroids using arithmetic means.                                                                                                                \\

            Empty-cluster handling       &
            Reinitializes empty clusters during the clustering process.                                                                                                      \\

            Distortion computation       &
            Computes clustering distortion, including a streaming variant for memory-safe distance evaluation.                                                               \\

            Scikit-learn-style interface &
            Provides methods such as \varname{fit()}, \varname{predict()}, \varname{transform()}, \varname{fit_predict()}, \varname{fit_transform()}, and \varname{score()}. \\
            \bottomrule
      \end{tabularx}
\end{table}

The clustering module supports several registered divergence functions, as shown in Table~\ref{tab:supported_bregman_divergences}.

\begin{table}[H]
      \centering
      \caption{Supported Bregman Divergences}
      \label{tab:supported_bregman_divergences}
      \renewcommand{\arraystretch}{1.2}
      \begin{tabularx}{\textwidth}{
                  >{\raggedright\arraybackslash}p{3.5cm}
                  >{\raggedright\arraybackslash}X
            }
            \toprule
            \textbf{Identifier} & \textbf{Divergence Function} \\
            \midrule
            \textbf{euclidean}  &
            Squared Euclidean divergence.                      \\

            \textbf{gkl}        &
            Generalized Kullback--Leibler divergence.          \\

            \textbf{is}         &
            Itakura--Saito divergence.                         \\

            \textbf{logistic}   &
            Logistic loss-based divergence.                    \\
            \bottomrule
      \end{tabularx}
\end{table}

Before clustering is performed, the \textbf{validate\_divergence\_domain()} function checks whether the input matrix contains finite values and satisfies the mathematical domain requirements of the selected divergence. This validation step helps prevent invalid clustering operations, especially for divergences that require non-negative or strictly positive input values.

\subsection{Machine Learning Module}

\hs The machine learning module supports the F-step of the \KFCProcedure{} framework. Its main role is to provide local predictive models that can be trained separately within the clusters produced by the K-step. The module includes abstract model interfaces, native local models, scikit-learn adapters, and a registry mechanism for resolving models by name. Table~\ref{tab:machine_learning_module_components} summarizes the main components of this module.

\begin{table}[H]
      \centering
      \caption{Machine Learning Module Components}
      \label{tab:machine_learning_module_components}
      \renewcommand{\arraystretch}{1.2}
      \begin{tabularx}{\textwidth}{
                  >{\raggedright\arraybackslash}p{3.5cm}
                  >{\raggedright\arraybackslash}X
            }
            \toprule
            \textbf{Component} & \textbf{Description}                                                             \\
            \midrule
            BaseLocalModel     &
            Defines the abstract interface used by local models in the F-step.                                    \\

            MeanRegressor      &
            Provides a simple native local regression model based on a mean prediction strategy.                  \\

            SklearnLocalModel  &
            Wraps compatible scikit-learn estimators so they can be used as local models inside the KFC pipeline. \\

            LocalModelFactory  &
            Provides a registry mechanism for resolving local models by name.                                     \\
            \bottomrule
      \end{tabularx}
\end{table}

The function \varname{register_all_sklearn_models()} automatically registers compatible scikit-learn classifiers and regressors using standardized snake-case identifiers. Examples include \varname{linear_regression}, \varname{logistic_regression}, \varname{random_forest_regressor}, and \varname{k_neighbors_classifier}. This registration mechanism allows users to select local models through configuration names instead of manually instantiating estimator objects.

\subsection{KFC Workflow Steps}

\hs The KFC pipeline is organized into three main workflow steps: clustering, local model training, and prediction aggregation. These steps correspond to the K-step, F-step, and C-step of the framework. Table~\ref{tab:kfc_workflow_steps} summarizes the role of each step in the pipeline.

\begin{table}[H]
      \centering
      \caption{KFC Workflow Steps}
      \label{tab:kfc_workflow_steps}
      \renewcommand{\arraystretch}{1.2}
      \begin{tabularx}{\textwidth}{
                  >{\raggedright\arraybackslash}p{2.5cm}
                  >{\raggedright\arraybackslash}p{3.3cm}
                  >{\raggedright\arraybackslash}X
            }
            \toprule
            \textbf{Step}                   & \textbf{Main Function} & \textbf{Description}                                                                                                                  \\
            \midrule
            KStep                           &
            Divergence-based clustering     &
            Receives divergence identifiers or divergence objects, fits one Bregman K-Means model for each divergence, and stores the fitted clustering models and training cluster assignments.             \\

            FStep                           &
            Cluster-specific model training &
            Uses the cluster labels generated by KStep. For each divergence and observed cluster, it resolves a local model, trains it on the corresponding data subset, and stores the fitted local models. \\

            CStep                           &
            Prediction aggregation          &
            Receives the prediction matrix generated by FStep, resolves the selected combiner, fits the combiner when required, and produces the final prediction through aggregation.                       \\
            \bottomrule
      \end{tabularx}
\end{table}

During prediction, \textbf{FStep.predict()} returns a prediction matrix of shape \textbf{(n\_samples, n\_divergences)}. Each column represents one divergence-specific predictive view. The cluster assignments are used internally to route each sample to the appropriate local model, but the final prediction matrix contains one column per divergence rather than one column per cluster.

\subsection{Combiner Module}

\hs The combiner module is responsible for the C-step of the \KFCProcedure{} framework. Its main function is to combine the prediction outputs produced by the cluster-specific models in the F-step. The module supports both regression and classification tasks through separate registered combiner strategies. \tabref{tab:registered_prediction_combiners} summarizes the registered prediction combiners used in the framework.

\begin{table}[H]
      \centering
      \caption{Registered Prediction Combiners}
      \label{tab:registered_prediction_combiners}
      \resizebox{\textwidth}{!}{%
            \begin{tabularx}{\textwidth}{
                        >{\raggedright\arraybackslash}p{2.5cm}
                        >{\raggedright\arraybackslash}p{3.3cm}
                        >{\raggedright\arraybackslash}X
                  }
                  \toprule
                  \textbf{Task}  & \textbf{Combiner Identifier}  & \textbf{Description}                                                    \\
                  \midrule
                  Regression     & \textbf{mean}                 & Computes the row-wise arithmetic mean of prediction columns.            \\
                  Regression     & \textbf{weighted\_mean}       & Combines prediction columns using linear-regression-based weights.      \\
                  Regression     & \textbf{stacking\_regressor}  & Applies a meta-regression model to combine prediction outputs.          \\
                  Regression     & \textbf{gradientcobra}        & Uses \textbf{GradientCOBRA} as a COBRA-based regression combiner.       \\
                  Regression     & \textbf{mixcobra}             & Uses \textbf{MixCOBRARegressor} as a mixed-space regression combiner.   \\
                  \midrule
                  Classification & \textbf{majority\_vote}       & Applies hard voting across prediction columns for classification.       \\
                  Classification & \textbf{stacking\_classifier} & Applies logistic-regression-based stacking for classification.          \\
                  Classification & \textbf{combined\_classifier} & Uses \textbf{CombinedClassifier} as a classification aggregation model. \\
                  \bottomrule
            \end{tabularx}%
      }
\end{table}

For COBRA-based KFC combiners, the corresponding COBRA model is called with \textbf{as\_predictions=True}. This setting is necessary because the combiner does not receive the original feature matrix directly. Instead, it receives the prediction matrix generated by the F-step. Therefore, the COBRA-based combiner operates in prediction space when it is used inside the KFC pipeline.

\subsection{Subsystem Summary}

\hs The KFC subsystem implements the clusterwise learning part of the framework through explicit packages for clustering, local model wrapping, pipeline steps, and prediction combiners. Its implementation is modular: divergence selection, local model choice, and final aggregation strategy can be changed independently through factory registries.

% ============================

\section{COBRA Aggregation Subsystem}
\label{sec:cobra_subsystem}

\hs The COBRA aggregation subsystem implements prediction-space aggregation models and reusable aggregation components. Unlike the KFC subsystem, which first partitions the input space into clusters, the COBRA subsystem primarily works by constructing or receiving a prediction matrix and then aggregating calibration targets according to distance-based kernel weights. This makes the subsystem suitable for combining the outputs of multiple predictive models in a flexible and modular way.

\subsection{Subsystem Architecture}

\hs The subsystem is organized around several model-level components that support regression, classification, and stacking-style aggregation. Table~\ref{tab:cobra_model_components} summarizes the main components of the COBRA aggregation subsystem.

\begin{table}[H]
      \centering
      \caption{Main Components of the COBRA Aggregation Subsystem}
      \label{tab:cobra_model_components}
      \begin{tabularx}{\textwidth}{lX}
            \toprule
            \textbf{Component}              & \textbf{Role in the Subsystem}                                             \\
            \midrule
            \textbf{GradientCOBRA}          &
            Provides regression aggregation based on prediction-space distances and kernel-weighted calibration targets. \\

            \textbf{MixCOBRARegressor}      &
            Combines input-space and prediction-space distance information for regression aggregation.                   \\

            \textbf{CombinedClassifier}     &
            Provides classification aggregation by combining the predictions of multiple classification models.          \\

            \textbf{CombinedClassifierFast} &
            Provides a faster variant of the combined classification aggregation procedure.                              \\
            \bottomrule
      \end{tabularx}
\end{table}

In addition to these model-level components, the subsystem includes reusable infrastructure for estimator management, data splitting, normalization, distance computation, kernel functions, loss functions, optimization, aggregation, and shared type definitions. These reusable components reduce code duplication and make the aggregation models easier to extend and maintain.


\subsection{Core Engine Structure}

\hs The COBRA core engine is designed as a factory-driven component system. It separates the main aggregation workflow into reusable components for distance computation, kernel weighting, loss evaluation, optimization, estimator management, and data splitting. This modular design makes the subsystem easier to extend and allows different aggregation strategies to reuse the same underlying infrastructure. Table~\ref{tab:cobra_core_engine_components} summarizes the main components of the COBRA core engine.

\begin{table}[H]
      \centering
      \caption{Core Components of the COBRA Engine}
      \label{tab:cobra_core_engine_components}
      \begin{tabularx}{\textwidth}{lX}
            \toprule
            \textbf{Component}        & \textbf{Role in the Core Engine}                                                                                                                                                                                                                \\
            \midrule
            \textbf{Adapters}         &
            Provide kernel-adapter mechanisms. \textbf{OneParameterKernelAdapter} scales one distance matrix using a bandwidth parameter, while \textbf{TwoParameterKernelAdapter} combines two distance matrices using coefficients $\alpha$ and $\beta$.                              \\

            \textbf{Aggregators}      &
            Combine calibration targets using computed weights. \textbf{WeightedMeanAggregator} is used for regression aggregation, while \textbf{WeightedVoteAggregator} is used for classification aggregation.                                                                       \\

            \textbf{Cross-validation} &
            Provides fold-generation strategies, including \textbf{KFoldCV}, \textbf{StratifiedKFoldCV}, and \textbf{TimeSeriesCV}. These components support internal validation and parameter-selection procedures.                                                                    \\

            \textbf{Distances}        &
            Implements distance measures such as Euclidean, Manhattan, Minkowski, Cosine, and Hamming distances. Hamming distance is especially useful for comparing classification prediction vectors.                                                                                 \\

            \textbf{Estimators}       &
            Provides estimator abstraction and registration through components such as \textbf{BaseEstimator}, \textbf{MeanRegressor}, \textbf{SklearnEstimator}, and \textbf{EstimatorFactory}. This allows scikit-learn estimators to be used consistently within the COBRA workflow. \\

            \textbf{Kernels}          &
            Defines kernel functions used to convert distances into aggregation weights. The available kernels include radial/RBF, exponential, reverse-cosh, Cauchy, Epanechnikov, biweight, triweight, triangular, COBRA threshold, and naive kernels.                                \\

            \textbf{Losses}           &
            Provides objective functions for model evaluation and optimization, including MSE, MAE, Huber, quantile, log loss, and hinge loss.                                                                                                                                          \\

            \textbf{Normalizers}      &
            Provides reusable normalization components, including standard normalization and min--max normalization. Shared preprocessing utilities are also used to compute scalar normalization constants where required.                                                             \\

            \textbf{Optimizers}       &
            Tunes aggregation parameters such as bandwidth, $\alpha$, and $\beta$. The optimizer layer includes \textbf{GridSearchOptimizer}, \textbf{GradientDescentOptimizer}, \textbf{MomentumOptimizer}, and \textbf{AdamOptimizer}.                                                \\

            \textbf{Splitters}        &
            Creates training and calibration partitions through strategies such as \textbf{RandomHoldoutSplitter} and \textbf{OverlapSplitter}.                                                                                                                                         \\
            \bottomrule
      \end{tabularx}
\end{table}


\subsection{GradientCOBRA Implementation}

\hs\textbf{GradientCOBRA} is implemented as a scikit-learn-compatible regressor for prediction-space aggregation. Its default estimator pool includes linear regression, ridge regression with cross-validation, lasso regression with cross-validation, k-nearest-neighbor regression, random forest regression, and support vector regression. These base estimators are used to generate prediction outputs that are later combined through the COBRA aggregation mechanism.

The \textbf{fit()} method supports two execution modes:

\begin{itemize}
      \item \textbf{Raw-feature mode:} The model receives the original feature matrix, splits the data into estimator-training and calibration subsets, trains the base estimators, and constructs a prediction matrix on the calibration data.

      \item \textbf{Prediction-matrix mode:} When \textbf{as\_predictions=True}, the input is treated as an existing prediction matrix. In this case, base estimator training is skipped, and the model directly performs aggregation on the provided predictions.
\end{itemize}

After the prediction matrix has been prepared, \textbf{GradientCOBRA} computes a scalar normalization constant, constructs a distance matrix in prediction space, applies a one-parameter kernel adapter, evaluates kernel-based aggregation weights, and optimizes the bandwidth parameter through cross-validation. This design allows the model to operate both as a standalone regression aggregator and as a combiner inside the \KFCProcedure{} pipeline.


\subsection{MixCOBRARegressor Implementation}

\hs\textbf{MixCOBRARegressor} extends the prediction-space aggregation idea by using both input-space and prediction-space distances. Its default estimator pool includes linear regression, ridge, lasso, k-nearest-neighbor regression, random forest regression, and support vector regression.

The class supports two modes controlled by \textbf{one\_parameter}:

\begin{itemize}
      \item \textbf{One-parameter mode:} Input features and prediction features are concatenated after normalization, and a single bandwidth-like parameter is optimized.

      \item \textbf{Two-parameter mode:} Input-space distance and prediction-space distance are computed separately and combined by the two-parameter adapter using $\alpha$ and $\beta$.
\end{itemize}

Therefore, \textbf{MixCOBRARegressor} should be described as a mixed-space distance aggregation method, not as a method that combines multiple kernel strategies.

\subsection{CombinedClassifier Implementation}

\hs\textbf{CombinedClassifier} implements a COBRA-style classifier. Its default estimator pool includes logistic regression, decision tree classifier, support vector classifier, and k-nearest-neighbor classifier. The default distance is Hamming distance, the default kernel is RBF, and the default aggregator is weighted vote.

The class supports both raw-feature mode and \textbf{as\_predictions=True} mode. In raw-feature mode, it trains base classifiers and constructs prediction vectors. In prediction-matrix mode, the given input matrix is used directly as the classifier prediction space.

The file also contains \textbf{CombinedClassifierFast}, an optional extension that can build a FAISS index when FAISS is available and \textbf{use\_faiss=True}. This is an acceleration path for nearest-neighbor search, not the default classifier path.

\subsection{Subsystem Summary}

\hs The COBRA subsystem implements a modular prediction-space aggregation framework. Its main classes reuse core components for splitting, estimator management, normalization, distance computation, kernel weighting, optimization, loss evaluation, and aggregation. The same subsystem can be used independently or as a KFC combiner through the C-step.


% ============================

\section{Orchestration Layer}
\label{sec:orchestration_layer}

\hs The orchestration layer provides the high-level estimator classes used to coordinate training and prediction in the framework. Its main role is to connect the K-step, F-step, and C-step into a unified workflow, while also allowing direct COBRA-based aggregation models to be used when required.

\subsection{Main Classes}

\hs The orchestration layer exposes three main classes, as summarized in Table~\ref{tab:orchestration_main_classes}.

\begin{table}[H]
      \centering
      \caption{Main Classes in the Orchestration Layer}
      \label{tab:orchestration_main_classes}
      \renewcommand{\arraystretch}{1.2}
      \begin{tabularx}{\textwidth}{
                  >{\raggedright\arraybackslash}p{3.2cm}
                  >{\raggedright\arraybackslash}p{3.0cm}
                  >{\raggedright\arraybackslash}X
            }
            \toprule
            \textbf{Class}               & \textbf{Task} & \textbf{Description}                             \\
            \midrule
            \textbf{KFCProcedure}        &
            Regression or Classification &
            General estimator with a configurable \textbf{task} parameter.                                  \\

            \textbf{KFCRegressor}        &
            Regression                   &
            Convenience wrapper that initializes \textbf{KFCProcedure} with \textbf{task="regression"}.     \\

            \textbf{KFCClassifier}       &
            Classification               &
            Convenience wrapper that initializes \textbf{KFCProcedure} with \textbf{task="classification"}. \\
            \bottomrule
      \end{tabularx}
\end{table}

\subsection{KFCProcedure Workflow}

\hs The \KFCProcedure{} workflow consists of two main phases: training and prediction. During training, the data are separated into two subsets. The first subset is used for divergence-based clustering and local model training, while the second subset is used to calibrate the final combiner. During prediction, new samples are first assigned to divergence-specific clusters, routed through the corresponding local models, and then aggregated by the fitted C-step combiner.

Figure~\ref{fig:kfc_training_prediction_workflow} summarizes the training and prediction workflow of \KFCProcedure{}.

\begin{figure}[H]
      \centering
      \begin{tikzpicture}[
                  font=\small,
                  node distance=0.55cm,
                  box/.style={
                              draw,
                              rounded corners,
                              align=center,
                              text width=5.7cm,
                              minimum height=0.75cm
                        },
                  title/.style={
                              font=\bfseries,
                              align=center
                        },
                  arrow/.style={-Latex, thick}
            ]

            % Training column
            \node[title] (train_title) at (0,0) {Training Workflow};
            \node[box, below=0.35cm of train_title] (t1) {Input training data $(X, y)$};
            \node[box, below=of t1] (t2) {Convert inputs into NumPy arrays};
            \node[box, below=of t2] (t3) {Split data into K/F training subset and C-step calibration subset};
            \node[box, below=of t3] (t4) {Fit \textbf{KStep} and obtain divergence-specific cluster assignments};
            \node[box, below=of t4] (t5) {Fit \textbf{FStep} using cluster-specific local models};
            \node[box, below=of t5] (t6) {Generate calibration prediction matrix};
            \node[box, below=of t6] (t7) {Fit \textbf{CStep} using the prediction matrix and calibration targets};
            \node[box, below=of t7] (t8) {Trained \KFCProcedure{} model};

            \draw[arrow] (t1) -- (t2);
            \draw[arrow] (t2) -- (t3);
            \draw[arrow] (t3) -- (t4);
            \draw[arrow] (t4) -- (t5);
            \draw[arrow] (t5) -- (t6);
            \draw[arrow] (t6) -- (t7);
            \draw[arrow] (t7) -- (t8);

            % Prediction column
            \node[title] (pred_title) at (7.0,0) {Prediction Workflow};
            \node[box, below=0.35cm of pred_title] (p1) {Input new data $X$};
            \node[box, below=of p1] (p2) {Check that K-step, F-step, and C-step are fitted};
            \node[box, below=of p2] (p3) {Assign divergence-specific clusters using \textbf{KStep}};
            \node[box, below=of p3] (p4) {Route samples through cluster-specific local models using \textbf{FStep}};
            \node[box, below=of p4] (p5) {Construct divergence-level prediction matrix};
            \node[box, below=of p5] (p6) {Apply fitted \textbf{CStep} combiner};
            \node[box, below=of p6] (p7) {Final prediction output};

            \draw[arrow] (p1) -- (p2);
            \draw[arrow] (p2) -- (p3);
            \draw[arrow] (p3) -- (p4);
            \draw[arrow] (p4) -- (p5);
            \draw[arrow] (p5) -- (p6);
            \draw[arrow] (p6) -- (p7);

      \end{tikzpicture}
      \caption{Training and prediction workflow of \KFCProcedure{}}
      \label{fig:kfc_training_prediction_workflow}
\end{figure}

The C-step may use different types of combiners, including simple rule-based combiners, stacking-based combiners, and COBRA-based combiners. Therefore, the \KFCProcedure{} workflow should be understood as a general clusterwise learning pipeline, not as a pipeline that always ends with COBRA aggregation. COBRA aggregation is used only when the selected combiner corresponds to a COBRA-based option, such as \textbf{"gradientcobra"}, \textbf{"mixcobra"}, or \textbf{"combined\_classifier"}.

\subsection{Direct COBRA Aggregation Workflow}

\hs The direct COBRA aggregation workflow can be used independently from the KFC pipeline. In raw-feature mode, the COBRA model trains a pool of base estimators and constructs a prediction matrix from their outputs. In prediction-matrix mode, the model receives an existing prediction matrix directly and skips base estimator training.

Figure~\ref{fig:direct_cobra_workflow} summarizes the direct COBRA aggregation workflow.

\begin{figure}[H]
      \centering
      \begin{tikzpicture}[
                  font=\small,
                  node distance=0.7cm,
                  box/.style={
                              draw,
                              rounded corners,
                              align=center,
                              text width=7.8cm,
                              minimum height=0.8cm
                        },
                  smallbox/.style={
                              draw,
                              rounded corners,
                              align=center,
                              text width=5.2cm,
                              minimum height=0.8cm
                        },
                  arrow/.style={-Latex, thick}
            ]

            \node[smallbox] (raw) at (-3.1,0) {Raw-feature mode\\Input features and target};
            \node[smallbox] (pred) at (3.1,0) {Prediction-matrix mode\\Input existing prediction matrix};

            \node[box, below=1.1cm of raw, xshift=3.1cm] (matrix) {Construct or receive the prediction matrix};

            \node[box, below=of matrix] (dist) {Compute distances in prediction space\\or in both input and prediction spaces for \textbf{MixCOBRARegressor}};
            \node[box, below=of dist] (kernel) {Transform distances using kernel adapters and kernel functions};
            \node[box, below=of kernel] (optim) {Optimize aggregation parameters through validation};
            \node[box, below=of optim] (aggregate) {Aggregate calibration targets using distance-based weights};
            \node[box, below=of aggregate] (output) {Final COBRA prediction output};

            \draw[arrow] (raw) -- node[left, font=\scriptsize] {fit estimators} (matrix);
            \draw[arrow] (pred) -- node[right, font=\scriptsize] {\textbf{as\_predictions=True}} (matrix);
            \draw[arrow] (matrix) -- (dist);
            \draw[arrow] (dist) -- (kernel);
            \draw[arrow] (kernel) -- (optim);
            \draw[arrow] (optim) -- (aggregate);
            \draw[arrow] (aggregate) -- (output);

      \end{tikzpicture}
      \caption{Direct COBRA aggregation workflow}
      \label{fig:direct_cobra_workflow}
\end{figure}

Prediction-matrix mode is especially important for integration with the KFC pipeline. In this case, the F-step has already generated a prediction matrix, so the COBRA-based combiner can operate directly on these predictions instead of retraining a new pool of base estimators.


\subsection{Workflow Summary}

\hs The orchestration layer provides a unified entry point for KFC training and inference while also supporting direct COBRA-based aggregation workflows. The \KFCProcedure{} workflow creates divergence-level predictions through clustering and local predictive models, while the COBRA workflow aggregates prediction matrices using distance, kernel, optimization, loss, and aggregation components. Together, these workflows provide a flexible structure for comparing standard supervised learning, clusterwise learning, and COBRA-based aggregation within the same package.

% ============================

\section{Utility Layer}
\label{sec:utility_layer}

\hs The utility layer provides supporting functions used by the KFC and COBRA subsystems. These utilities assist with logging, component resolution, data splitting, normalization, estimator management, and distance computation. However, they do not implement the main learning algorithms of the framework.

\subsection{Top-Level KFC Utilities}

The top-level KFC utilities support execution control and configuration resolution within the KFC pipeline. Table~\ref{tab:kfc_top_level_utilities} summarizes the main utility components.

\begin{table}[H]
      \centering
      \caption{Top-Level KFC Utility Components}
      \label{tab:kfc_top_level_utilities}
      \begin{tabularx}{\textwidth}{
                  >{\raggedright\arraybackslash}p{3.5cm}
                  >{\raggedright\arraybackslash}X
            }
            \toprule
            \textbf{Utility Component}    & \textbf{Description}                                                                                                              \\
            \midrule
            \textbf{Logger}               &
            Provides lightweight execution messages with verbosity levels. It supports informational, debugging, and tracing messages, as well as timed execution monitoring. \\

            \textbf{Resolution Utilities} &
            Resolve Bregman clustering configurations into clustering objects and construct K-step configurations from user-defined parameters.                               \\
            \bottomrule
      \end{tabularx}
\end{table}

The current top-level utility layer does not provide a general preprocessing pipeline for missing-value handling, feature scaling, or categorical encoding. These preprocessing steps are handled externally in the experimental workflow or by the user before the framework is called.

\subsection{COBRA Utilities}

The COBRA utilities support the construction and execution of COBRA-based aggregation workflows. They are mainly used to prepare training and calibration contexts, fit estimator pools, compute prediction matrices, normalize prediction spaces, and resolve factory-based components. Table~\ref{tab:cobra_utility_components} summarizes the main COBRA utility components.

\begin{table}[H]
      \centering
      \caption{COBRA Utility Components}
      \label{tab:cobra_utility_components}
      \begin{tabularx}{\textwidth}{
                  >{\raggedright\arraybackslash}p{3.5cm}
                  >{\raggedright\arraybackslash}X
            }
            \toprule
            \textbf{Utility Component}       & \textbf{Description}                                                                                                                        \\
            \midrule
            \textbf{Preprocessing Utilities} &
            Provide helper functions for overlapping data splits, normalization constant computation, estimator-name formatting, and conversion of optimization history into tabular form. \\

            \textbf{Resolution Utilities}    &
            Resolve estimators, kernels, splitters, losses, distances, aggregators, and training contexts. They also support fitting and prediction for estimator pools.                   \\

            \textbf{Distance Utilities}      &
            Provide optimized distance-computation support, including a Numba-accelerated Hamming distance matrix function for efficient comparison of prediction vectors.                 \\
            \bottomrule
      \end{tabularx}
\end{table}


\section{Testing and Validation}
\label{sec:testing_validation}

\hs The project includes an automated test suite under \textbf{tests/}. In the current implementation, the automated tests mainly cover the COBRA core components. Separate automated tests for the \KFCProcedure{} subsystem and the full end-to-end \KFCProcedure{} workflow are not yet included.

\subsection{Test Suite Structure}

\hs The test suite is organized as shown in Figure~\ref{fig:cobra_test_suite_structure}.

\begin{figure}[H]
      \centering
      \begin{minipage}{0.85\textwidth}
            \begin{Verbatim}[fontsize=\footnotesize, baselinestretch=0.9]
                  tests/
                  └── cobra/
                  ├── test_p01_splitters.py
                  ├── test_p02_estimators.py
                  ├── test_p03_normalizers.py
                  ├── test_p04_distances.py
                  ├── test_p05_kernel_adapters.py
                  ├── test_p06_kernels.py
                  ├── test_p07_losses.py
                  ├── test_p08_cv.py
                  ├── test_p09_optimizers.py
                  └── test_p10_aggregators.py
            \end{Verbatim}
      \end{minipage}
      \caption{Test suite structure for COBRA components}
      \label{fig:cobra_test_suite_structure}
\end{figure}

Figure~\ref{fig:cobra_test_suite_structure} shows that the current automated test suite is organized around the COBRA module. The tests are separated by component type, including splitters, estimators, normalizers, distances, kernel adapters, kernels, losses, cross-validation, optimizers, and aggregators.

\subsection{Covered Components}

\hs The automated tests cover several COBRA components, as summarized in Table~\ref{tab:covered_cobra_test_components}.

\begin{table}[H]
      \centering
      \caption{COBRA Components Covered by Automated Tests}
      \label{tab:covered_cobra_test_components}
      \renewcommand{\arraystretch}{1.2}
      \begin{tabularx}{\textwidth}{
                  >{\raggedright\arraybackslash}p{3.2cm}
                  >{\raggedright\arraybackslash}X
            }
            \toprule
            \textbf{Component}        & \textbf{Test Coverage}                                                                                                  \\
            \midrule
            \textbf{Splitters}        &
            Holdout splitting, overlap splitting, reproducibility, factory creation, and parameter validation.                                                  \\

            \textbf{Estimators}       &
            Mean regressor behavior, scikit-learn wrapper behavior, factory creation, and interface contract checks.                                            \\

            \textbf{Normalizers}      &
            Standard normalization, min--max normalization, constant-feature handling, and factory resolution.                                                  \\

            \textbf{Distances}        &
            Euclidean, Manhattan, Minkowski, Cosine, and Hamming distance behavior.                                                                             \\

            \textbf{Kernel Adapters}  &
            One-parameter scaling, two-parameter distance combination, parameter updates, and input validation.                                                 \\

            \textbf{Kernels}          &
            Factory aliases and behavior of radial, exponential, reverse-cosh, Cauchy, Epanechnikov, biweight, triweight, triangular, naive, and COBRA kernels. \\

            \textbf{Losses}           &
            MSE, MAE, Huber, log loss, hinge loss, and quantile loss.                                                                                           \\

            \textbf{Cross-validation} &
            K-fold counts, disjoint validation sets, full sample coverage, leakage prevention between training and validation sets, and reproducibility.        \\

            \textbf{Optimizers}       &
            Gradient descent, momentum, Adam, grid search, optimizer factories, and optimization history.                                                       \\

            \textbf{Aggregators}      &
            Weighted mean, weighted vote, fallback behavior, NaN handling, batch aggregation, and empty-value checks.                                           \\
            \bottomrule
      \end{tabularx}

\end{table}

Table~\ref{tab:covered_cobra_test_components} shows that the current automated tests focus on the internal behavior of COBRA-related components. These tests help verify that the main building blocks used in the aggregation layer behave as expected.

\subsection{Current Testing Gaps}

\hs Although the automated tests cover several COBRA core components, some parts of the full framework are not yet covered by separate automated tests. Table~\ref{tab:current_testing_gaps} summarizes the main testing gaps in the current repository.

\begin{table}[H]
      \centering
      \caption{Current Testing Gaps}
      \label{tab:current_testing_gaps}
      \renewcommand{\arraystretch}{1.2}
      \begin{tabularx}{\textwidth}{
                  >{\raggedright\arraybackslash}p{4.2cm}
                  >{\raggedright\arraybackslash}X
            }
            \toprule
            \textbf{Scope}                            & \textbf{Missing Test Coverage}                                                 \\
            \midrule
            BregmanKMeans, divergence, and clustering &
            Separate automated tests for Bregman K-Means behavior and clustering under each supported divergence are not yet included. \\
            KStep, FStep, and CStep                   &
            The individual KFC pipeline steps are not yet covered by independent component-level tests.                                  \\

            Full KFCProcedure workflow                &
            End-to-end tests for complete KFCProcedure training and prediction are not yet included.                                     \\

            Regression and classification wrappers    &
            Dedicated automated tests for KFCRegressor and KFCClassifier are not yet included.                                           \\

            Notebook execution                        &
            Experimental notebooks are not yet executed as part of an automated validation pipeline.                                     \\
            \bottomrule
      \end{tabularx}

\end{table}

Table~\ref{tab:current_testing_gaps} indicates that the current testing coverage should be described as partial rather than complete. The repository includes automated tests for several COBRA core components, but it should not be presented as having full system-level testing unless additional end-to-end KFC tests are added in future work.

\subsection{Validation Through Notebooks}

\hs The project also uses experiment notebooks as a practical validation mechanism. These notebooks demonstrate regression and classification experiments using baseline models, COBRA-based models, and KFC configurations. However, notebook execution is not currently integrated into the automated pytest suite or continuous integration workflow.

\subsection{Testing Summary}

\hs The implemented automated tests provide coverage for the COBRA core component layer, as shown in Figure~\ref{fig:cobra_test_suite_structure} and summarized in Table~\ref{tab:covered_cobra_test_components}. The KFC subsystem and full orchestration workflow are demonstrated in notebooks, but they are not yet covered by automated unit or integration tests in the repository, as summarized in Table~\ref{tab:current_testing_gaps}.

\section{Documentation and CI/CD Pipeline}
\label{sec:documentation_cicd}

\hs This section describes the documentation and automation infrastructure included in the project. The project provides README-level documentation, inline docstrings, package metadata, pytest configuration, and a GitHub Actions workflow for testing, building, and publishing the package.

\subsection{Documentation Assets}

\hs The project documentation is mainly provided through project-level documentation, inline source documentation, package metadata, and experiment notebooks. Table~\ref{tab:documentation_assets} summarizes the verified documentation assets.

\begin{table}[H]
      \centering
      \caption{Documentation Assets}
      \label{tab:documentation_assets}
      \renewcommand{\arraystretch}{1.2}
      \begin{tabularx}{\textwidth}{
                  >{\raggedright\arraybackslash}p{3.8cm}
                  >{\raggedright\arraybackslash}X
            }
            \toprule
            \textbf{Asset}       & \textbf{Description}                                                                                        \\
            \midrule
            README documentation &
            Provides a project overview, feature list, installation instructions, quick-start examples, and high-level component descriptions. \\

            Inline docstrings    &
            Provide explanatory documentation for several modules and classes, including parameters, workflows, and design intent.             \\

            Package metadata     &
            Defines package information, dependencies, optional extras, build configuration, and pytest settings.                              \\

            Experiment notebooks &
            Provide executable examples and experimental usage patterns for regression and classification tasks.                               \\
            \bottomrule
      \end{tabularx}
\end{table}

No separate Sphinx, MkDocs, ReadTheDocs, or GitHub Pages documentation configuration is currently included. Therefore, automated documentation generation should not be claimed unless such tooling is added in future development.

\subsection{Package Configuration}

\hs The package configuration defines the project metadata, dependency requirements, optional extras, build settings, and testing configuration. Table~\ref{tab:package_configuration_summary} summarizes the verified package configuration.

\begin{table}[H]
      \centering
      \caption{Package Configuration Summary}
      \label{tab:package_configuration_summary}
      \renewcommand{\arraystretch}{1.2}
      \begin{tabularx}{\textwidth}{
                  >{\raggedright\arraybackslash}p{4.0cm}
                  >{\raggedright\arraybackslash}X
            }
            \toprule
            \textbf{Configuration Item} & \textbf{Verified Description}                                                                          \\
            \midrule
            Package name                &
            The package is configured under the name \textbf{kfc-procedure}.                                                                     \\

            Python requirement          &
            The package requires Python version \textbf{3.11} or higher.                                                                         \\

            Core dependencies           &
            The package includes core scientific and machine learning dependencies such as NumPy, pandas, scikit-learn, matplotlib, and XGBoost. \\

            Optional extras             &
            Optional dependency groups are provided for core usage, COBRA functionality, development tools, and full installation.               \\

            Test configuration          &
            The pytest configuration defines the test directory used for automated test discovery.                                               \\
            \bottomrule
      \end{tabularx}
\end{table}

\subsection{Continuous Integration and Publishing Workflow}

\hs The project includes a GitHub Actions workflow for package testing, building, and publishing. The workflow is triggered by Git version tags that follow a semantic versioning pattern. Table~\ref{tab:cicd_workflow_steps} summarizes the verified workflow steps.

\begin{table}[H]
      \centering
      \caption{CI/CD Workflow Steps}
      \label{tab:cicd_workflow_steps}
      \renewcommand{\arraystretch}{1.2}
      \begin{tabularx}{\textwidth}{
                  >{\raggedright\arraybackslash}p{1.6cm}
                  >{\raggedright\arraybackslash}X
            }
            \toprule
            \textbf{Step} & \textbf{Description}                                             \\
            \midrule
            1             &
            Check out the repository source code.                                            \\

            2             &
            Set up the Python execution environment.                                         \\

            3             &
            Install the package with development and COBRA-related dependencies.             \\

            4             &
            Run the automated test suite using \textbf{pytest}.                              \\

            5             &
            Build the package distribution.                                                  \\

            6             &
            Upload the built distribution to PyPI using \textbf{twine} and a PyPI API token. \\
            \bottomrule
      \end{tabularx}
\end{table}

This workflow can be described as a tag-triggered test, build, and publishing pipeline. It is not currently configured to run on every push or pull request, and it does not include verified linting, static analysis, or documentation deployment steps.

\subsection{Automation Tools}

\hs The verified automation tools are summarized in Table~\ref{tab:automation_tools}.

\begin{table}[H]
      \centering
      \caption{Automation Tools Used in the Project}
      \label{tab:automation_tools}
      \renewcommand{\arraystretch}{1.2}
      \begin{tabularx}{\textwidth}{
                  >{\raggedright\arraybackslash}p{3.2cm}
                  >{\raggedright\arraybackslash}X
            }
            \toprule
            \textbf{Tool}  & \textbf{Role}                                                        \\
            \midrule
            GitHub Actions &
            Provides the tag-triggered automation workflow for testing, building, and publishing. \\

            pytest         &
            Runs the automated test suite.                                                        \\

            setuptools     &
            Supports package configuration and packaging through the project metadata.            \\

            build          &
            Generates the package distribution files.                                             \\

            twine          &
            Uploads the built distribution to PyPI.                                               \\
            \bottomrule
      \end{tabularx}
\end{table}

Although a pylint configuration is present, the current GitHub Actions workflow does not run pylint, flake8, ruff, or another linter. Therefore, linting and static code quality checks should not be presented as verified parts of the current CI/CD pipeline.

\subsection{Summary}

\hs The project includes functional package configuration and a tag-triggered CI/CD workflow for testing, building, and publishing the package. Documentation is currently provided through README-level documentation, inline docstrings, experiment notebooks, and package metadata. More advanced documentation deployment, broader CI checks, and automated linting would require additional configuration in future development.

\section{Chapter Summary}
\label{sec:chapter_summary}

\hs This chapter presented the implementation of the \textbf{kfc\_procedure} framework as a modular Python research library. The package combines two related subsystems: the \KFCProcedure{} clusterwise learning subsystem and the \GradientCOBRA{} prediction-space aggregation subsystem.

The KFC subsystem provides the main clusterwise learning pipeline. It includes Bregman K-Means clustering, registered Bregman divergences, local model wrappers, K-step/F-step/C-step workflow classes, and task-specific combiner strategies. The COBRA subsystem provides complementary aggregation models, including \textbf{GradientCOBRA}, \textbf{MixCOBRARegressor}, \textbf{CombinedClassifier}, and reusable components for splitting, estimator management, distance computation, kernel weighting, optimization, normalization, loss evaluation, cross-validation, and aggregation.

The orchestration layer coordinates the KFC workflow through \textbf{KFCProcedure}, \textbf{KFCRegressor}, and \textbf{KFCClassifier}. During training, the implementation separates the data into a K/F training subset and a C-step calibration subset, fits divergence-specific clustering models, trains cluster-specific local models, constructs a divergence-level prediction matrix, and fits the selected combiner. During prediction, new samples are assigned to clusters, routed through local models, and aggregated by the fitted C-step.

The implementation applies modular software engineering principles through factory registries, component-level abstractions, scikit-learn-style interfaces, composition-based workflow construction, strategy-based aggregation selection, and separation between clustering, local modeling, aggregation, optimization, and utility functions. The current automated test suite focuses mainly on COBRA core components, while full KFC subsystem and end-to-end workflow tests remain areas for future extension.

Finally, the project includes package metadata, README-level documentation, experiment notebooks, and a tag-triggered GitHub Actions workflow for testing, building, and publishing the package. Overall, the implementation translates the methodology from Chapter~\ref{ch:methodology} into a reusable and extensible Python framework, while leaving room for further improvements in automated testing, documentation, and continuous integration coverage.

